\section{Contoh: Menelusuri Satu Bab dengan Alur OBE}
\label{sec:contoh-bab1}

Agar alur belajar pada Bagian~\ref{sec:cara-menggunakan} menjadi
konkret, berikut contoh penerapannya pada satu bab, yaitu
Bab~\ref{chap:register} (Register 8086 dan Register Flag) yang
menempati minggu ke-3.

\subsection{Langkah 1: Membaca Sub-CPMK}

Bab tersebut menetapkan \SubCPMK{1.3} dengan indikator: mengidentifikasi
fungsi register umum, register segmen, register \eng{pointer} dan
\eng{index}, serta register \eng{flag} beserta pengaruhnya terhadap
eksekusi program. Dari rumusan ini mahasiswa sudah tahu bahwa yang
akan diuji bukan sekadar kemampuan menghafal nama register, melainkan
kemampuan \emph{memprediksi} perilaku program ketika register berubah.

\subsection{Langkah 2: Membaca Materi dan Menandai Pertanyaan}

Saat membaca, tandai hal yang belum jelas. Contoh pertanyaan yang
wajar muncul pada bab tersebut:

\begin{itemize}
    \item Mengapa \reg{AX} disebut akumulator, dan apakah ia wajib
          dipakai pada operasi aritmatika?
    \item Apa bedanya \reg{SP} dan \reg{BP}? Keduanya menunjuk stack,
          tetapi mengapa perlu dua register?
    \item Jika \flag{CF} dan \flag{OF} sama-sama menunjukkan
          ``kelebihan'', mengapa keduanya dibedakan?
\end{itemize}

Pertanyaan seperti ini dibawa ke kelas. Inilah inti
\eng{flipped classroom}: waktu tatap muka dipakai untuk menjawab
kebingungan, bukan untuk membacakan ulang isi buku.

\subsection{Langkah 3: Menjalankan Contoh Program}

Bagian Contoh Program pada bab tersebut meminta mahasiswa menjalankan
program penjumlahan dan memprediksi nilai flag sebelum program
dijalankan. Tabel~\ref{tab:contoh-prediksi} memperlihatkan bentuk
tabel pengamatan yang diisi setelah program dijalankan.

\begin{table}[H]
    \centering
    \caption{Contoh tabel prediksi dan hasil pengamatan \eng{flag}
    untuk \asm{ADD AL, BL}.}
    \label{tab:contoh-prediksi}
    \small
    \begin{tabular}{@{}cccccc@{}}
        \toprule
        \textbf{AL} & \textbf{BL} & \textbf{Hasil} &
        \textbf{Prediksi ZF} & \textbf{Prediksi CF} &
        \textbf{Keterangan} \\
        \midrule
        \hex{7F} & \hex{01} & \hex{80} & 0 & 0 &
        \textit{overflow} bertanda, \flag{OF} = 1 \\
        \addlinespace
        \hex{FF} & \hex{01} & \hex{00} & 1 & 1 &
        \eng{carry} keluar bit ke-8 \\
        \addlinespace
        \hex{40} & \hex{40} & \hex{80} & 0 & 0 &
        hasil bertanda negatif, \flag{SF} = 1 \\
        \addlinespace
        \hex{10} & \hex{10} & \hex{20} & 0 & 0 &
        hasil biasa, semua flag utama 0 \\
        \bottomrule
    \end{tabular}
\end{table}

\subsection{Langkah 4: Mengisi Jurnal dan Checklist}

Setelah praktikum, mahasiswa mengisi jurnal praktikum sesuai format
Tabel~\ref{tab:format-jurnal}, lalu menandai butir Checklist
Pencapaian Kompetensi pada akhir bab. Butir yang belum bertanda
menjadi daftar tugas belajar sebelum kuis.

\subsection{Langkah 5: Menghubungkan ke Asesmen}

Bagian Evaluasi pada bab tersebut memuat kisi-kisi yang menautkan
setiap soal ke indikator \SubCPMK{1.3}. Dengan membaca kisi-kisi itu
sebelum mengerjakan, mahasiswa mengetahui bobot relatif antara
kemampuan menghafal (\dec{15}\%) dan kemampuan menganalisis
(\dec{85}\%). Informasi ini membantu mengarahkan usaha belajar pada
kemampuan yang benar-benar dinilai.

\begin{infobox}
Perhatikan pola yang berulang: dimulai dari \emph{outcome}, diakhiri
dengan \emph{asesmen}, dan seluruh kegiatan di antaranya adalah
sarana. Pola inilah yang membedakan belajar dalam kerangka OBE dari
belajar menghafal materi.
\end{infobox}
